\documentclass[10pt,a4paper]{amsart}
\usepackage[margin=19mm]{geometry}
\usepackage{amsmath,amssymb,mathtools}
\usepackage[scr=rsfs,cal=euler]{mathalfa}
\usepackage{xfrac}
\usepackage[unicode,hidelinks,bookmarks=false]{hyperref}
\newcommand{\ie}{i.e.\ }
\newcommand{\cf}{cf.\ }
\newcommand{\resp}{respectively\ }
\newcommand{\Subsection}{\subsection}
\providecommand{\nobreakdash}{\nobreak\hbox{-}}
\setlength{\emergencystretch}{2em}
\multlinegap0pt
\usepackage{bm}
\usepackage{bbm}
\usepackage{dsfont}
\usepackage{xspace}
\usepackage{cancel}
\usepackage{mathtools}
\usepackage{amsthm}
\makeatletter
\@ifundefined{theorem}{\newtheorem{theorem}{Theorem}}{}
\@ifundefined{remark}{\newtheorem{remark}[theorem]{Remark}}{}
\makeatother
\def \dcl {\operatorname{dcl}}
\def \tp {tp}
\title[On definable Galois theory and definable Galois cohomology i]{On definable Galois theory and definable Galois cohomology in the totally transcendental setting}
\author{David Meretzky}
\date{Source: arXiv:2609.17828v1 (CC BY 4.0)}
\begin{document}
\maketitle

\noindent\emph{Source and licence.} On definable Galois theory and definable Galois cohomology in the totally transcendental setting. Version arXiv:2609.17828v1 (CC BY 4.0), distributed under the Creative Commons Attribution 4.0 International licence, as recorded in the arXiv licence metadata for this exact version and shown on the abstract page https://arxiv.org/abs/2609.17828v1. Full rights notice: licenses/arxiv-2609.17828-CC-BY-4.0.txt.\par\smallskip

\noindent\textit{Editorial scope.} Counted unit: the Remark whose source label is \texttt{None} in the article (source0.tex lines 444--446, source label \texttt{None}), delivered together with the article's own complete original proof of it (source0.tex lines 448--463, one \texttt{proof} environment of 16 lines). No mathematics is written, repaired or re-derived: statement and proof are the article's own text line for line. Required context retained uncounted: none. Every definition, display and label of those lines is retained unchanged. Omitted from this package and not redistributed: 1--443, 447--447, 464--991. Nothing inside a retained excerpt is omitted or altered: every retained body is delivered exactly as the article's lines are written, so no omission receipt of any kind accompanies this package. Citations: the retained excerpts carry no citation command at all, so no bibliography entry of the article is transcribed and no reference is redistributed. Reference adaptation: none was needed and none was made; every reference inside the delivered unit resolves inside it, so no unresolved reference or citation is produced. Printed slips kept literal: no printed word, symbol, hypothesis or number of the counted statement or of its proof was altered. Layout and class: the article's own class is \texttt{amsart} with packages including amssymb, hyperref, url, xcolor, tikz-cd; the wrapper is the lane's pinned \texttt{amsart} editorial wrapper at 10pt on A4, which restates the theorem-like environments the retained text needs and adds the article's own macro definitions that the retained text needs (\texttt{\textbackslash dcl}, \texttt{\textbackslash tp}), with extra wrapper packages (bm, bbm, dsfont, xspace, cancel, mathtools). The delivered numbering is the unit's own. Assets: no retained span contains a figure environment, a graphics-inclusion command, a PDF-page inclusion, a table or a TikZ picture; no image file is copied into this package, the package declares no asset dependency and no image file is redistributed with it. Licence: version arXiv:2609.17828v1 of this article is distributed under the Creative Commons Attribution 4.0 International licence, as recorded in the arXiv licence metadata for this exact version and shown on the abstract page https://arxiv.org/abs/2609.17828v1. A full rights notice is delivered at licenses/arxiv-2609.17828-CC-BY-4.0.txt. Characters: every retained excerpt is pure ASCII, so the delivered engine, XeLaTeX with its default Latin Modern text font, reports no missing character.
\begin{remark}
The tuple $\bar{c}$ in most cases will come from $X$. In the case that $X$ and $A$ are both empty, the assumption that the language has multiple constants guarantees that $dcl(X,A)$ is still at least countably infinite: Note that if there are at least two constants in the language then finite binary sequences may be coded.
\end{remark}

\begin{proof}
Note that $A$ and $B$ are small.  By assumption, the collection of $A\cup B$-formulas $$\{\psi_Y(\bar{z})\} \cup \{\neg \exists \bar{x}_{\bar{a}} \in X^{|\bar{x}_{\bar{a}}|}(\varphi_{\bar{a}}(\bar{b},\bar{x}_{\bar{a}},\bar{z})) : \bar{b} \in B^{\bar{y}_{\bar{a}}}\}_{\bar{a} \in Y}$$ is inconsistent. Therefore it is finitely inconsistent by compactness. Thus there are finitely many $\bar{a}_1,...,\bar{a}_m \in Y$ such that $$\models \forall \bar{z} \left( \psi_Y(\bar{z}) \rightarrow \bigvee_{i=1}^m \ \exists \bar{x}_{\bar{a}_i} \in X^{|\bar{x}_{\bar{a}_i}|} \  \varphi_{\bar{a}_i}(\bar{b}_{\bar{a}_i},\bar{x}_{\bar{a}_i},\bar{z})\right).$$ 

Let $e_1,...,e_m$ be distinct elements from $dcl(X,A)$ (which exist by the preceding remark) associated to variables $u_1,...,u_m$. Then for each $\bar{a} \in Y$, there exists $\bar{c}_{\bar{a}} \in X^{|\bar{x}_{\bar{a}_i}|}$ and $e_j \in dcl(X,A)$ such that $$\models\bigvee_{i=1}^m (\varphi_{\bar{a}_i}(\bar{b}_{\bar{a}_i},\bar{c}_{\bar{a}},\bar{a}) \wedge e_i=e_j)$$ and $\bar{a}$ is the only such solution.

Denote the above formula by $\varphi(\bar{y},\bar{x},\bar{z})$ where
we let $\bar{b} =(b_{\bar{a}_1},...,b_{\bar{a}_m})$ in the variables $\bar{y} = (\bar{y}_{\bar{a}_1},...,\bar{y}_{\bar{a}_m})$ and $\bar{x} = (\bar{x}_{\bar{a}_1},u_1,....,\bar{x}_{\bar{a}_m},u_m, u)$ with $u$ being the variable filled by the element $e_j$. It is clear that $\varphi(\bar{b}, \bar{x},\bar{z})$ defines the graph of a definable function.

Let $\bar{a} \in Y$ now and $\bar{c} \in dcl(X,A)$ such that $\models \varphi(\bar{b},\bar{c},\bar{a})$. Then $\varphi(\bar{y},\bar{c},\bar{a}) \in \tp(\bar{b}/dcl(X,A)\cup Y)$. As $T$ is totally transcendental, by definability of types over arbitrary subsets of $\bar{M}$, $\varphi(\bar{b},\bar{x},\bar{z})$ is definable over $dcl(X,A) \cup Y$. 

We now make some replacements in notation and we are finished. 

The set defined by $\varphi(\bar{b},\bar{x},\bar{z})$, is defined by a $dcl(X,A) \cup Y$-formula which we will now write as $\varphi_0(\bar{b},\bar{c}, \bar{x},\bar{z})$. Where $\bar{b}=(b_1,...,b_n)$ is a new tuple of parameters from $Y$, $\bar{c}$, in say variables $\bar{v}$, is a new tuple of parameters from $\dcl(A,X)$ and $\varphi_0(\bar{y},\bar{v},\bar{x}, \bar{z})$ is $A$-definable. As before, $\varphi_0(\bar{b},\bar{c},\bar{x}, \bar{z})$ is the graph of a definable function. Now both $\bar{v}$ and $\bar{x}$ are sorts on $\dcl(A,X)$ so rewrite them as just $\bar{x}$. Then $\varphi_0(\bar{b},\bar{x},\bar{z})$ is still clearly a the graph of a definable function $f(\bar{b},\bar{x}) = \bar{z}$.

The tuple $\bar{x} = (\bar{x}_{\bar{a}_1},u_1,....,\bar{x}_{\bar{a}_m},u_m,u,\bar{v})$. So the $\bar{x}$ variable ranges over an $A$-definable subset $X_0$ of $\dcl(A,X)$.
\end{proof}

\end{document}
